% =============================================================================
% Manual do Morphe -- RASCUNHO (23/09/2026)
%
% Capitulos 1 a 3 escritos, em capitulos/. Os demais ainda trazem, numa caixa
% cinza, o que vao conter. Figuras em TikZ, em figuras/.
%
% Compila no Overleaf com pdfLaTeX (padrao). Arquivo principal: main.tex.
% =============================================================================
\documentclass[11pt,a4paper,oneside]{report}

\usepackage[utf8]{inputenc}
\usepackage[T1]{fontenc}
\usepackage{lmodern}
\usepackage[brazil]{babel}
\usepackage[a4paper,margin=2.5cm]{geometry}
\usepackage{microtype}
\usepackage{amsmath}
% nomes de arquivo nao hifenizam: folga para a linha nao estourar a margem
\setlength{\emergencystretch}{3em}
\usepackage{xcolor}
\usepackage{booktabs}
\usepackage{tabularx}
\usepackage{array}
\usepackage{dirtree}
\usepackage[most]{tcolorbox}
\usepackage{tikz}
\usetikzlibrary{positioning,arrows.meta,fit,backgrounds,calc,
                decorations.pathreplacing,shapes.geometric,matrix}
\usepackage[hidelinks]{hyperref}

% Cores e estilos comuns a todas as figuras.
% Uma cor por camada, a mesma em todo o manual:
%   cliente (PC) = azul, servidor (HPS) = verde, FPGA = laranja.
\definecolor{cCliente}{HTML}{2F6DB5}
\definecolor{cHps}{HTML}{2E8B57}
\definecolor{cFpga}{HTML}{D2691E}
\definecolor{cRede}{HTML}{6B6B6B}
\definecolor{cFuturo}{HTML}{9A9A9A}
% Uma cor por operacao, quando a figura separa os aceleradores entre si.
\definecolor{opConv}{HTML}{C0392B}
\definecolor{opFft}{HTML}{7D3C98}
\definecolor{opIir}{HTML}{B7950B}

\tikzset{
  >={Stealth[length=2.2mm]},
  every picture/.append style={font=\small},
  caixa/.style={draw, rounded corners=2pt, align=center, inner sep=4pt,
                minimum height=8mm},
  cliente/.style={caixa, draw=cCliente, fill=cCliente!10},
  hps/.style={caixa, draw=cHps, fill=cHps!10},
  fpga/.style={caixa, draw=cFpga, fill=cFpga!12},
  rede/.style={caixa, draw=cRede, fill=cRede!8},
  mem/.style={draw=cFpga!80!black, fill=white, minimum height=6mm,
              inner sep=2pt, font=\scriptsize\ttfamily},
  futuro/.style={caixa, draw=cFuturo, dashed, fill=white, text=cFuturo!70!black},
  grupo/.style={draw, rounded corners=4pt, inner sep=6pt},
  seta/.style={->, thick},
  rotulo/.style={font=\scriptsize, align=center},
}


% Caixa "o que esta secao vai conter" -- sai do documento final.
\newtcolorbox{previsto}{colback=black!4,colframe=black!25,boxrule=0.4pt,
  arc=1pt,left=6pt,right=6pt,top=4pt,bottom=4pt,fontupper=\small\itshape,
  before upper={\textbf{\upshape Conteúdo previsto.}\ }}

% Nome de arquivo ou comando, sempre em pé (dentro das caixas em itálico,
% o itálico da fonte monoespaçada juntaria "--" num travessão).
\newcommand{\arq}[1]{{\upshape\ttfamily\detokenize{#1}}}
% Opcao de linha de comando: \opt{check} escreve --check (hifens separados).
\newcommand{\opt}[1]{{\upshape\ttfamily -{}-\detokenize{#1}}}

\title{\textbf{Morphe}\\[4pt]
  \large Plataforma de Processamento Digital de Sinais em FPGA\\
  Manual do sistema: funcionamento e estrutura\\[10pt]
  \normalsize\textcolor{red!70!black}{Rascunho: capítulos 1 a 3 escritos; os demais, em estrutura}}
\author{Antonio Nicassio Santos Lima\\
  \small Estágio no LabPS/UEFS --- supervisão: Prof.\ Armando S.\ Sanca}
\date{23 de setembro de 2026}

\begin{document}
\maketitle

\begin{abstract}
\begin{previsto}
Uma página: o que é o Morphe (operações de PDS executadas na FPGA de uma
DE1-SoC e comparadas com a referência em software), a quem o manual se destina
(professor, monitor e aluno que vai estender o sistema), o que foi herdado do
TCC de Carlos Valadão e o que o estágio acrescentou. Os números que resumem a
plataforma: cinco operações em hardware, sinais de qualquer tamanho por
blocos, até \(\sim\)92\,dB de SNR na FFT, seis clientes simultâneos por placa
sem perda de resposta.
\end{previsto}
\end{abstract}

\tableofcontents
\listoffigures

% Capitulos ja escritos: um arquivo por capitulo.
% =============================================================================
\chapter{Introdução}
\label{cap:introducao}
% =============================================================================

\section{O que é o Morphe}

O Morphe é uma plataforma para executar operações de Processamento Digital de
Sinais (PDS) em hardware e comparar o resultado com o que o computador calcula.
O aluno prepara um sinal numa janela do aplicativo, escolhe a operação e recebe
de volta duas respostas lado a lado: a que a FPGA de uma placa DE1-SoC calculou
em ponto fixo e a que o próprio computador calculou em ponto flutuante. A
diferença entre as duas é o objeto de estudo.

São cinco as operações executadas em hardware:

\begin{description}
  \item[Convolução] linear de dois sinais de até 1024 amostras cada, com saída
    de até 2047 amostras.
  \item[Filtro FIR] a mesma convolução, com a resposta ao impulso vinda de um
    projetista de filtros por especificação.
  \item[Filtro IIR] cascata de até 16 seções de segunda ordem (ordem até 32),
    também com projetista próprio.
  \item[FFT] transformada de Fourier discreta de 1024 pontos.
  \item[IFFT] a transformada inversa, no mesmo núcleo da FFT.
\end{description}

O hardware trabalha com blocos de tamanho fixo (1024 amostras), mas o
aplicativo aceita sinais de qualquer tamanho: divide o sinal em blocos, manda
cada um à placa e recompõe o resultado, com técnicas que dão a mesma resposta
que uma operação única daria (capítulo~\ref{cap:cliente}).

\section{Para que serve no ensino de PDS}

Nas disciplinas de PDS as operações são estudadas com números reais, e o
computador as calcula em ponto flutuante de 64 bits, que para quase todo fim
didático é o mesmo que calcular com números reais. Um circuito digital dedicado
não trabalha assim: usa ponto fixo, com um número pequeno de bits para a parte
fracionária, e cada multiplicação precisa ser arredondada ou truncada.

O Morphe torna essa diferença visível e mensurável. Com ele o aluno pode:

\begin{itemize}
  \item ver o erro de quantização de uma FFT em decibéis e observar que ele
    melhora 6\,dB a cada bit de faixa aproveitada (capítulo~\ref{cap:numerica});
  \item provocar uma saturação ou um estouro de acumulador e reconhecer o
    sintoma na saída;
  \item comparar um filtro IIR projetado em ponto flutuante com o mesmo filtro
    executado em ponto fixo, e entender por que o hardware usa seções de segunda
    ordem e não um polinômio de ordem alta;
  \item medir quanto tempo cada operação leva na placa e relacionar esse tempo
    com o número de operações do algoritmo.
\end{itemize}

\section{Histórico}

O Morphe nasceu como Trabalho de Conclusão de Curso de Carlos Valadão na UEFS,
\emph{Implementação de Ferramentas de Computação Numérica para o Processamento
Digital de Sinais em Hardware Reconfigurável}. A arquitetura, o núcleo de
convolução \texttt{conv1d}, a integração do núcleo de FFT, o servidor em C e o
cliente gráfico em Python são dele, e o projeto é distribuído sob a licença GNU
GPL versão~3.

A partir de setembro de 2026 a plataforma foi continuada num estágio curricular
no LabPS/UEFS, por Antonio Nicassio Santos Lima, sob supervisão do Prof.\ Armando
S.\ Sanca. O objetivo do estágio foi tornar a plataforma reproduzível a partir do
repositório e utilizável por uma turma de alunos que não conhece os detalhes de
baixo nível. A Tabela~\ref{tab:historico} resume o que mudou; o registro dia a
dia está em \arq{docs/DIARIO.md} e as correções feitas no projeto original, com
os motivos, em \arq{docs/RESSALVAS.md}.

\begin{table}[htbp]
  \centering
  \small
  \caption{O que o estágio acrescentou ao projeto original.}
  \label{tab:historico}
  \begin{tabularx}{\textwidth}{@{}lXX@{}}
    \toprule
    & \textbf{Projeto original} & \textbf{Hoje} \\
    \midrule
    Tamanho & 128 amostras na convolução e no FIR & 1024 amostras em todas as
      operações, e sinais maiores por blocos \\
    Operações & convolução, FIR, FFT & as três, mais IFFT e IIR \\
    Filtros & sem projetista por especificação & projetistas FIR e IIR por especificação \\
    Precisão & entrada da FFT sem normalização & normalização da entrada: de
      50 para 92\,dB de SNR na FFT \\
    Preparar a placa & passos manuais, conhecidos pelo autor & um comando,
      \arq{morphe-up.sh} \\
    Laboratório & uma estação, uma placa & duas placas, vários clientes ao mesmo
      tempo, instalação para a turma e Quartus em todos os PCs \\
    Reprodutibilidade & o que funcionava estava só na estação & tudo sai do
      repositório, com manifesto de integridade \\
    \bottomrule
  \end{tabularx}
\end{table}

\section{Visão geral}

A Figura~\ref{fig:visao-geral} é a figura de referência deste manual: todo o
resto detalha uma das partes dela.

\begin{figure}[htbp]
  \centering
  \begin{tikzpicture}[node distance=5mm]
  % Computadores
  \node[cliente, minimum width=30mm] (pc1) {PC do laboratório\\\scriptsize aplicativo Morphe};
  \node[cliente, minimum width=30mm, below=of pc1] (pc2) {PC do laboratório\\\scriptsize aplicativo Morphe};
  \node[cliente, minimum width=30mm, below=of pc2] (est) {Estação\\\scriptsize aplicativo + \texttt{morphe-up.sh}};

  % Rede
  \node[rede, right=14mm of pc2, minimum height=34mm, text width=18mm] (sw) {Rede do\\laboratório};

  % Placas
  \node[hps, right=16mm of sw.north east, anchor=north west, yshift=4mm, text width=24mm]
        (hps1) {HPS\\\scriptsize ARM + Linux\\\texttt{morphe\_server}};
  \node[fpga, right=10mm of hps1, text width=18mm, minimum height=16mm] (fp1) {FPGA\\\scriptsize aceleradores};
  \node[hps, right=16mm of sw.south east, anchor=south west, yshift=-4mm, text width=24mm]
        (hps2) {HPS\\\scriptsize ARM + Linux\\\texttt{morphe\_server}};
  \node[fpga, right=10mm of hps2, text width=18mm, minimum height=16mm] (fp2) {FPGA\\\scriptsize aceleradores};

  \begin{scope}[on background layer]
    \node[grupo, fit=(hps1)(fp1), label={[font=\scriptsize\bfseries]above:DE1-SoC (placa 1)}] (b1) {};
    \node[grupo, fit=(hps2)(fp2), label={[font=\scriptsize\bfseries]above:DE1-SoC (placa 2)}] (b2) {};
  \end{scope}

  \draw[<->, thick] (hps1) -- node[rotulo, above]{AXI} (fp1);
  \draw[<->, thick] (hps2) -- node[rotulo, above]{AXI} (fp2);

  % Rede
  \foreach \p in {pc1, pc2, est}
    \draw[<->, thick, cRede] (\p.east) -- (\p.east -| sw.west);
  \draw[<->, thick, cRede] (sw.east |- hps1) -- node[rotulo, above]{TCP 5000} (hps1.west);
  \draw[<->, thick, cRede] (sw.east |- hps2) -- node[rotulo, below]{TCP 5000\\SSH} (hps2.west);

  % JTAG
  \draw[->, thick, dashed, cFpga]
    (est.south) -- ++(0,-6mm) -| node[rotulo, pos=0.25, below]{USB-Blaster (JTAG): programa a FPGA e mantém o \emph{tether}} (fp2.south);
\end{tikzpicture}

  \caption{Visão geral: os computadores do laboratório falam com as placas pela
    rede (TCP, porta 5000); a estação programa a FPGA pelo USB-Blaster e sobe o
    servidor por SSH.}
  \label{fig:visao-geral}
\end{figure}

Cada \textbf{PC do laboratório} roda o aplicativo Morphe, escrito em Python. Ele
não precisa de nenhum software da Intel: fala com as placas só pela rede.

Cada \textbf{placa DE1-SoC} tem, no mesmo chip, um processador ARM rodando Linux
(o HPS, \emph{Hard Processor System}) e uma FPGA. No processador roda o
\arq{morphe_server}, que recebe as requisições pela rede na porta 5000; na FPGA
estão os aceleradores que fazem as contas. Os dois lados conversam por pontes
internas ao chip (seção~\ref{sec:pontes}).

A \textbf{estação} é o computador ligado às placas pelo cabo USB-Blaster. É dela
que se roda o \arq{morphe-up.sh}, que programa a FPGA pelo JTAG, compila e
reinicia o servidor na placa por SSH e registra a placa para os clientes. A
estação também roda o aplicativo, como qualquer outro PC.

\section{Como ler este manual}

Quem só vai \textbf{usar} a plataforma precisa do capítulo~\ref{cap:cliente}
(o aplicativo) e do capítulo~\ref{cap:laboratorio} (preparar as placas). Quem vai
\textbf{estender} a plataforma --- acrescentar uma operação, mudar um acelerador,
recompilar o hardware --- deve ler os capítulos em ordem: a
arquitetura~(\ref{cap:arquitetura}), o hardware~(\ref{cap:hardware}), o
servidor~(\ref{cap:servidor}), o protocolo~(\ref{cap:protocolo}) e a
representação numérica~(\ref{cap:numerica}) são o que se precisa saber antes de
mexer em qualquer um deles.

Convenções usadas no texto:

\begin{itemize}
  \item nomes de arquivos, comandos e sinais do hardware aparecem em
    \arq{fonte monoespaçada}; caminhos são relativos à raiz do repositório;
  \item endereços e constantes do hardware aparecem em hexadecimal, com o
    prefixo \texttt{0x};
  \item \(\mathrm{Q}m.n\) é um número em ponto fixo com sinal, \(m\) bits de
    parte inteira e \(n\) de parte fracionária (capítulo~\ref{cap:numerica});
  \item todo número medido vem com a condição em que foi medido. Um número sem
    condição é de especificação, não de medida.
\end{itemize}

% =============================================================================
\chapter{Arquitetura em três camadas}
\label{cap:arquitetura}
% =============================================================================

\section{As camadas e o que cabe a cada uma}

O Morphe se divide em três camadas, uma por máquina ou por tipo de circuito
(Figura~\ref{fig:camadas}). Cada uma é escrita numa linguagem e tem uma
responsabilidade que as outras não repetem.

\begin{figure}[htbp]
  \centering
  % As tres camadas e os modulos de cada uma.
\begin{tikzpicture}[
    mod/.style={caixa, minimum width=27mm, minimum height=9mm, font=\scriptsize},
    faixa/.style={draw, rounded corners=4pt, inner sep=4mm},
  ]
  % --- Cliente -----------------------------------------------------------
  \node[mod, cliente] (jan)  at (0,0)     {janelas\\\texttt{*\_window.py}};
  \node[mod, cliente] (blo)  at (30mm,0)  {\texttt{blocos.py}\\sinais longos};
  \node[mod, cliente] (dsp)  at (60mm,0)  {\texttt{dsp\_core.py}\\ponto fixo e referência};
  \node[mod, cliente] (proj) at (90mm,0)  {projetistas\\\texttt{fir/iir\_design}};
  \node[mod, cliente] (pro)  at (120mm,0) {\texttt{morphe\_protocol}\\cliente TCP};

  % --- Servidor ----------------------------------------------------------
  \node[mod, hps, minimum width=52mm] (srv) at (60mm,-24mm)
        {\texttt{morphe\_server.c}\\valida, copia, dispara, espera, lê};
  \node[mod, hps] (mk) at (0,-24mm) {marca\\\texttt{fpga-preparada}};
  \draw[seta] (srv.west) -- node[rotulo, above]{consulta} (mk.east);

  % --- FPGA --------------------------------------------------------------
  \node[mod, fpga] (conv) at (15mm,-48mm)  {\texttt{conv1d}\\convolução};
  \node[mod, fpga] (fir)  at (45mm,-48mm)  {\texttt{conv1d}\\FIR};
  \node[mod, fpga] (fft)  at (75mm,-48mm)  {IP FFT II\\FFT e IFFT};
  \node[mod, fpga] (iir)  at (105mm,-48mm) {\texttt{iir\_cascade}\\IIR};

  % Faixas, todas da mesma largura
  \begin{scope}[on background layer]
    \node[faixa, draw=cCliente, fill=cCliente!4, fit=(jan)(pro),
          label={[anchor=south west, inner sep=1pt, font=\scriptsize\bfseries, text=cCliente]north west:{PC --- Python}}] {};
    \node[faixa, draw=cHps, fill=cHps!4, fit=(mk)(srv)(pro.east |- srv),
          label={[anchor=south west, inner sep=1pt, font=\scriptsize\bfseries, text=cHps]north west:{HPS --- C, Linux}}] {};
    \node[faixa, draw=cFpga, fill=cFpga!4, fit=(conv)(iir)(jan.west |- conv)(pro.east |- iir),
          label={[anchor=south west, inner sep=1pt, font=\scriptsize\bfseries, text=cFpga]north west:{FPGA --- Verilog}}] {};
  \end{scope}

  \draw[seta] (jan) -- (blo);
  \draw[seta] (blo) -- (dsp);
  \draw[seta] (proj) -- (dsp);
  \draw[seta] (dsp.north) to[bend left=25] (pro.north);
  \draw[<->, thick, cRede] (pro.south) |- node[rotulo, pos=0.3, left]{TCP\\5000} (srv.east);
  \foreach \m in {conv, fir, fft, iir}
    \draw[seta] (srv.south) -- ++(0,-5mm) -| (\m.north);
  \node[rotulo, right] at ($(srv.south)+(1mm,-2.5mm)$) {\texttt{/dev/mem}};
\end{tikzpicture}

  \caption{As três camadas e os módulos de cada uma.}
  \label{fig:camadas}
\end{figure}

\begin{table}[htbp]
  \centering
  \small
  \caption{Responsabilidades de cada camada.}
  \label{tab:camadas}
  \begin{tabularx}{\textwidth}{@{}lXX@{}}
    \toprule
    & \textbf{Faz} & \textbf{Não faz} \\
    \midrule
    \textbf{Cliente} \newline \footnotesize Python, no PC
      & gera ou lê o sinal; projeta filtros; converte para ponto fixo; divide
        sinais longos em blocos; calcula a referência em ponto flutuante;
        compara e mostra
      & não sabe endereços de memória nem detalhes da FPGA \\
    \addlinespace
    \textbf{Servidor} \newline \footnotesize C, no HPS da placa
      & valida a requisição; copia os dados para as memórias da FPGA; dispara o
        acelerador; espera; devolve o resultado
      & não converte formatos numéricos (com uma exceção na FFT, seção
        \ref{sec:acel-fft}); não guarda nada entre requisições \\
    \addlinespace
    \textbf{Aceleradores} \newline \footnotesize Verilog, na FPGA
      & fazem a conta, em ponto fixo, sobre blocos de tamanho fixo
      & não conhecem a rede nem o tamanho do sinal que o aluno pediu \\
    \bottomrule
  \end{tabularx}
\end{table}

\subsection*{Três princípios que atravessam as camadas}

\paragraph{O servidor não guarda estado.} Cada operação chega completa --- os
dados e todos os parâmetros --- e sai completa. Não há sessão, reserva de placa
nem memória da requisição anterior. É isso que permite que vários clientes usem
a mesma placa intercalando operações, e que um cliente possa mudar de placa de
uma operação para a outra sem nada a transferir.

\paragraph{O ponto fixo é decidido no cliente.} O cliente converte o sinal para
o formato do acelerador (Q15.16 ou Q15.8) antes de enviar, e converte o
resultado de volta. O servidor copia inteiros de 32 bits sem interpretá-los,
com uma única exceção: aplicar o expoente de bloco da FFT, que só ele conhece
(seção~\ref{sec:acel-fft}). Fora isso, toda decisão numérica --- normalizar a
entrada, onde saturar, que fator aplicar à IFFT --- mora num lugar só, em
Python, onde é fácil de ler e de testar.

\paragraph{O hardware tem tamanho fixo; o software se adapta.} Os aceleradores
foram sintetizados para 1024 amostras. O servidor completa com zeros o que o
cliente não mandou e devolve só as amostras que fazem sentido; o cliente, por
sua vez, divide sinais maiores em blocos de 1024. Nenhum dos dois lados exige
que o aluno pense em 1024.

\subsection*{Os números de configuração existem em dois lugares}

Os limites do hardware --- 1024 amostras, 16 seções de IIR, os bits
fracionários de cada formato --- estão escritos em \arq{C/morphe_config.h}, para
o servidor, e repetidos em \arq{Python/morphe_config.py}, para o cliente. Os
dois arquivos precisam mudar juntos. Os endereços das memórias e registradores,
ao contrário, não são escritos à mão: vêm do \arq{C/hps_0.h}, gerado a partir do
projeto do hardware (seção~\ref{sec:pontes}).

\section{O caminho de uma operação}
\label{sec:caminho}

A Figura~\ref{fig:sequencia} acompanha uma convolução do clique na janela até o
gráfico. As outras operações seguem o mesmo roteiro, mudando só as memórias e o
formato numérico.

\begin{figure}[htbp]
  \centering
  % Diagrama de sequencia de uma operacao (ex.: convolucao).
\begin{tikzpicture}[
    ator/.style={caixa, minimum width=26mm, font=\scriptsize},
    msg/.style={->, thick},
    nota/.style={font=\scriptsize, align=left, fill=white, inner sep=1pt},
  ]
  \def\xa{0}\def\xb{36mm}\def\xc{74mm}\def\xd{112mm}
  \def\fim{-112mm}

  \node[ator, cliente] (a) at (\xa,0) {Janela\\(\texttt{conv\_window})};
  \node[ator, cliente] (b) at (\xb,0) {\texttt{morphe\_protocol}};
  \node[ator, hps]     (c) at (\xc,0) {\texttt{morphe\_server}};
  \node[ator, fpga]    (d) at (\xd,0) {\texttt{conv1d}\\e memórias};

  \foreach \n in {a,b,c,d}
    \draw[densely dashed, black!40] (\n.south) -- (\n.south |- 0,\fim);

  % Atividade no servidor
  \fill[cHps!25] ($(\xc,-26mm)+(-1mm,0)$) rectangle ($(\xc,-100mm)+(1mm,0)$);

  \draw[msg] (\xa,-12mm) -- node[above, rotulo]{\(x[n],\ h[n]\) em ponto flutuante} (\xb,-12mm);
  \draw[msg] (\xb,-17mm) arc[start angle=90, end angle=-90, radius=2.5mm]
             node[right, nota, pos=0.5, xshift=1mm] {codifica em Q15.16};
  \draw[msg, cRede] (\xb,-26mm) -- node[above, rotulo]{\texttt{connect}; cabeçalho de 20\,B + \emph{payload}} (\xc,-26mm);
  \draw[msg] (\xc,-32mm) arc[start angle=90, end angle=-90, radius=2.5mm]
             node[right, nota, pos=0.5, xshift=1mm] {valida magic, versão,\\opcode, tamanhos, marca};
  \draw[msg] (\xc,-44mm) -- node[above, rotulo]{escreve \(x\) e \(h\) nas SRAMs} (\xd,-44mm);
  \draw[msg] (\xc,-52mm) -- node[above, rotulo]{\emph{start} = 1 (ponte leve)} (\xd,-52mm);
  \draw[msg, dashed] (\xd,-60mm) -- node[above, rotulo]{lê \emph{done} até 1 (limite 5\,s)} (\xc,-60mm);
  \node[nota, right] at (\xd,-57mm) {calcula\\\(\approx\)214\,ms};
  \draw[msg, dashed] (\xd,-70mm) -- node[above, rotulo]{lê \(y\) da SRAM de saída} (\xc,-70mm);
  \draw[msg, cRede] (\xc,-82mm) -- node[above, rotulo]{cabeçalho de resposta + \(y\); \texttt{close}} (\xb,-82mm);
  \draw[msg] (\xb,-88mm) arc[start angle=90, end angle=-90, radius=2.5mm]
             node[right, nota, pos=0.5, xshift=1mm] {decodifica, remove\\o \emph{padding}};
  \draw[msg] (\xb,-100mm) -- node[above, rotulo]{\(y\) da placa \(\times\) \(y\) de referência} (\xa,-100mm);
\end{tikzpicture}

  \caption{O caminho de uma operação, do clique na janela ao gráfico.}
  \label{fig:sequencia}
\end{figure}

\begin{enumerate}
  \item \textbf{A janela entrega os sinais.} \(x[n]\) e \(h[n]\) chegam ao
    módulo de protocolo como vetores em ponto flutuante, do tamanho que o aluno
    escolheu.

  \item \textbf{O cliente codifica.} Cada amostra vira um inteiro de 32 bits em
    Q15.16: \(q = \operatorname{round}(v \cdot 2^{16})\), saturado na faixa do
    formato. Se o sinal passa de 1024 amostras, a divisão em blocos acontece
    antes deste passo, e os passos seguintes se repetem por bloco.

  \item \textbf{Uma conexão por operação.} O cliente abre uma conexão TCP com a
    placa, envia um cabeçalho de 20 bytes (operação, tipo dos dados, tamanhos) e
    os dados logo em seguida (capítulo~\ref{cap:protocolo}).

  \item \textbf{O servidor valida.} Confere a assinatura e a versão do
    protocolo, a operação, os tamanhos contra os limites do hardware e se a FPGA
    foi preparada pelo \arq{morphe-up.sh} (capítulo~\ref{cap:servidor}). Qualquer
    falha vira uma resposta de erro com uma mensagem legível, e a FPGA não é
    tocada.

  \item \textbf{O servidor escreve nas memórias.} Copia \(x\) e \(h\) para as
    memórias de entrada do acelerador, através da ponte HPS\(\to\)FPGA,
    completando com zeros até 1024 amostras.

  \item \textbf{O servidor dispara.} Escreve 1 no registrador \texttt{start}
    do acelerador, através da ponte leve. O acelerador detecta a subida desse
    sinal e começa.

  \item \textbf{O acelerador calcula} e, ao terminar, leva o sinal
    \texttt{done} a 1. A convolução de 1024 por 1024 amostras leva cerca de
    214\,ms; a FFT de 1024 pontos, cerca de 21\,ms (tempos de ida e volta
    medidos pelo \arq{morphe_ping} em 15/09/2026, com a placa sem outros
    clientes).

  \item \textbf{O servidor espera.} Lê o \texttt{done} a cada 50\,\textmu s. Se ele
    não subir em 5\,s, a operação é abandonada com o erro de tempo esgotado ---
    o sintoma típico de um bitstream que parou de funcionar
    (seção~\ref{sec:aceleradores}).

  \item \textbf{O servidor lê e responde.} Copia as \(n_x + n_h - 1\) amostras
    úteis da memória de saída, devolve o \texttt{start} a 0 e envia o cabeçalho
    de resposta e os dados. Fecha a conexão.

  \item \textbf{O cliente decodifica e compara.} Converte de volta para ponto
    flutuante, descarta o que for enchimento de zeros e calcula a referência em
    ponto flutuante com o NumPy. A janela mostra as duas e o erro entre elas.
\end{enumerate}

Dois fatos decorrem desse roteiro e aparecem na prática.

O tempo de uma operação \textbf{não depende do tamanho do sinal pedido}: o
acelerador sempre processa os blocos completos de 1024 amostras, e os zeros de
enchimento custam o mesmo que amostras de verdade. Pela construção do
\texttt{conv1d}, uma convolução de 10 por 10 amostras custa o mesmo que uma de
1024 por 1024.

Como o servidor atende \textbf{uma conexão de cada vez}, duas pessoas na mesma
placa não recebem erro: a segunda espera a primeira terminar, e o tempo de cada
uma soma. O capítulo~\ref{cap:placas} trata disso.

\section{Onde está cada camada no repositório}

\begin{center}
\small
\begin{tabular}{@{}ll@{}}
  \toprule
  \textbf{Camada} & \textbf{Onde} \\
  \midrule
  Cliente      & \arq{Python/} (aplicativo e módulos); testes em \arq{Python/ferramentas/} \\
  Servidor     & \arq{C/} (\arq{morphe_server.c} e cabeçalhos) \\
  Aceleradores & \arq{Quartus/} (RTL, projeto do Platform Designer, bitstream) \\
  Operação     & \arq{morphe-up.sh}, \arq{instala-quartus.sh}, \arq{deploy.sh}, na raiz \\
  \bottomrule
\end{tabular}
\end{center}

\noindent A árvore completa está no apêndice~\ref{ap:repositorio}.

% =============================================================================
\chapter{O hardware}
\label{cap:hardware}
% =============================================================================

\section{A placa DE1-SoC}

A DE1-SoC, da Terasic, é construída em torno de um Cyclone~V SoC da Intel
(\texttt{5CSEMA5F31C6}). O ``SoC'' do nome quer dizer que o chip reúne duas
coisas bem diferentes:

\begin{description}
  \item[HPS] (\emph{Hard Processor System}): um processador ARM Cortex-A9 de
    dois núcleos, com memória e periféricos próprios, que roda um Linux gravado
    no cartão SD da placa. No Morphe ele cuida da rede e roda o servidor.
  \item[FPGA] lógica programável, configurada pelo bitstream que o Quartus gera.
    No Morphe ela contém os aceleradores e as memórias em que eles trabalham.
\end{description}

Os dois lados estão no mesmo chip e se comunicam por pontes internas, sem
passar por nenhum fio da placa. É essa proximidade que permite ao servidor
escrever mil amostras numa memória da FPGA como se escrevesse num vetor em C.

A lógica da FPGA roda com o relógio de 50\,MHz da placa (\texttt{CLOCK\_50}).
A análise de tempo da última compilação (15/09/2026) mede uma frequência máxima
de 60,99\,MHz para esse domínio: há folga de cerca de 20\,\%. A
Tabela~\ref{tab:recursos} mostra quanto do chip o projeto ocupa.

\begin{table}[htbp]
  \centering
  \small
  \caption{Recursos da FPGA ocupados pelo projeto (compilação de 15/09/2026).}
  \label{tab:recursos}
  \begin{tabular}{@{}lrrr@{}}
    \toprule
    \textbf{Recurso} & \textbf{Usado} & \textbf{Disponível} & \textbf{Ocupação} \\
    \midrule
    Elementos lógicos (ALMs)      & 11\,140 & 32\,070 & 35\,\% \\
    Blocos DSP (multiplicadores)  & 30      & 87      & 34\,\% \\
    Blocos de memória M10K        & 165     & 397     & 42\,\% \\
    \bottomrule
  \end{tabular}
\end{table}

O que sobra é bastante: cerca de 230 blocos M10K livres, pouco menos de
300\,KB, o suficiente para uma memória de captura de dezenas de milhares de
amostras quando a aquisição pelo ADC for implementada (seção~\ref{sec:adc}).

\subsection*{O bitstream e a licença da FFT}

O bitstream versionado é \arq{Quartus/output_files/soc_system_time_limited.sof},
compilado com o Quartus Prime Lite 20.1. O núcleo de FFT é um IP da Intel que a
edição Lite só permite usar em modo de avaliação: o bitstream funciona
\textbf{enquanto o programador JTAG continuar conectado} à placa --- o
\emph{tether} --- e para depois de cerca de uma hora sem ele. Isso vale para o
bitstream inteiro, não só para a FFT: sem o \emph{tether}, a convolução também
passa a esgotar o tempo. O \arq{morphe-up.sh} mantém o \emph{tether} vivo
(capítulo~\ref{cap:laboratorio}).

\section{Pontes e mapa de memória}
\label{sec:pontes}

O HPS enxerga a FPGA como faixas de endereços. O Morphe usa duas pontes
(Figura~\ref{fig:mapa-memoria}):

\begin{description}
  \item[Ponte HPS\(\to\)FPGA] (\texttt{h2f\_axi\_master}), a partir de
    \texttt{0xC000\,0000}. Larga e rápida, é por onde passam os dados: as
    memórias de entrada e de saída de cada acelerador.
  \item[Ponte leve] (\emph{lightweight}), a partir de \texttt{0xFF20\,0000}.
    Estreita, feita para registradores de controle: os PIOs de
    \texttt{start}, \texttt{done} e de parâmetros.
\end{description}

\begin{figure}[htbp]
  \centering
  % Mapa de memoria. Enderecos de C/hps_0.h e C/address_map_arm.h.
\begin{tikzpicture}[
    reg/.style={draw, minimum width=34mm, minimum height=5.2mm, inner sep=1pt,
                font=\scriptsize\ttfamily, anchor=north},
    end/.style={font=\scriptsize\ttfamily, anchor=east, text=black!70},
  ]
  % ---- Ponte HPS->FPGA ---------------------------------------------------
  \node[font=\small\bfseries, align=center] at (0,8mm) {Ponte HPS\(\to\)FPGA\\\normalfont\scriptsize base \texttt{0xC000\,0000}, 32 bits};
  \foreach \i/\nome/\off/\cor in {
      0/fir\_yn/0x10000/opConv,
      1/conv1d\_yn/0x12000/opConv,
      2/fir\_hn/0x14000/opConv,
      3/fir\_xn/0x15000/opConv,
      4/conv1d\_hn/0x16000/opConv,
      5/conv1d\_xn/0x17000/opConv,
      6/fft\_yn\_re/0x18000/opFft,
      7/fft\_yn\_imag/0x19000/opFft,
      8/fft\_xn\_re/0x1A000/opFft,
      9/fft\_xn\_imag/0x1B000/opFft,
      10/iir\_xn/0x1C000/opIir,
      11/iir\_yn/0x1D000/opIir,
      12/iir\_coef/0x1E000/opIir}
  {
    \node[reg, fill=\cor!12] (h\i) at (0,-\i*5.2mm) {\nome};
    \node[end] at (h\i.west) {\off\ };
  }

  % ---- Ponte leve ---------------------------------------------------------
  \begin{scope}[xshift=72mm]
  \node[font=\small\bfseries, align=center] at (0,8mm) {Ponte leve (LW)\\\normalfont\scriptsize base \texttt{0xFF20\,0000}, PIOs};
  \foreach \i/\nome/\off/\cor in {
      0/fir\_error/0x00/opConv,
      1/fir\_done/0x10/opConv,
      2/fir\_start/0x20/opConv,
      3/conv1d\_done/0x30/opConv,
      4/conv1d\_start/0x40/opConv,
      5/fft\_bfp\_exponent/0x50/opFft,
      6/fft\_wrapper\_done/0x60/opFft,
      7/fft\_inverse/0x70/opFft,
      8/fft\_wrapper\_start/0x80/opFft,
      9/iir\_start/0x90/opIir,
      10/iir\_done/0xA0/opIir,
      11/iir\_error/0xB0/opIir,
      12/iir\_nsecoes/0xC0/opIir}
  {
    \node[reg, fill=\cor!12] (l\i) at (0,-\i*5.2mm) {\nome};
    \node[end] at (l\i.west) {\off\ };
  }
  \end{scope}

  % Legenda
  \begin{scope}[yshift=-72mm, xshift=-10mm]
    \foreach \x/\c/\t in {0/opConv/convolução e FIR, 38/opFft/FFT e IFFT, 72/opIir/IIR}
      { \fill[\c!12, draw=black] (\x mm,0) rectangle ++(4mm,3mm);
        \node[anchor=west, font=\scriptsize] at (\x mm+5mm,1.5mm) {\t}; }
  \end{scope}
\end{tikzpicture}

  \caption{Mapa de memória visto pelo servidor: memórias de dados atrás da
    ponte HPS\(\to\)FPGA e registradores de controle atrás da ponte leve.
    Os endereços são deslocamentos a partir da base de cada ponte.}
  \label{fig:mapa-memoria}
\end{figure}

\paragraph{As memórias têm duas portas.} Cada uma das treze memórias de dados
é uma RAM interna da FPGA de porta dupla: uma porta está ligada à ponte, para o
HPS, e a outra ao acelerador. O servidor e o acelerador nunca disputam a mesma
porta; o que os separa no tempo é o protocolo de \texttt{start} e \texttt{done}
(seção~\ref{sec:aceleradores}). As memórias de entrada guardam 1024 palavras de
32 bits (4\,KB); as de saída da convolução e do FIR, 2047 (8\,KB); a de
coeficientes do IIR, 5 palavras por seção para até 16 seções.

\paragraph{Como o servidor chega lá.} O servidor abre \arq{/dev/mem} e mapeia
as duas faixas no próprio espaço de endereços com \texttt{mmap}; a partir daí,
escrever \texttt{g\_conv\_xn[i]} é escrever na memória da FPGA. Por isso ele
precisa rodar como \texttt{root}. O tamanho da faixa mapeada é calculado em
tempo de compilação a partir do \arq{hps_0.h}, e verificações estáticas no
código recusam compilar se alguma memória for menor do que os limites de
\arq{morphe_config.h} exigem.

\paragraph{O \arq{hps_0.h} é gerado, nunca editado.} Os endereços da
Figura~\ref{fig:mapa-memoria} são decididos no Platform Designer, a ferramenta
do Quartus que monta o sistema, e ficam registrados no arquivo
\arq{soc_system.sopcinfo}. O \arq{Quartus/gen_hps_header.py} lê esse arquivo e
escreve o \arq{C/hps_0.h}; com \opt{check}, apenas confere se o cabeçalho
versionado ainda corresponde ao projeto. Um servidor compilado com endereços
velhos escreve nos lugares errados sem dar erro nenhum; a conferência existe
para que isso não aconteça.

\medskip
\noindent\fbox{\parbox{\dimexpr\textwidth-2\fboxsep-2\fboxrule}{\small
\textbf{Cuidado com o mapa do ``DE1-SoC Computer''.} Os exemplos da Intel
para a placa usam \texttt{0xC800\,0000} como base da memória da FPGA. No
Morphe a base é \texttt{0xC000\,0000}. Código copiado desses exemplos lê e
escreve fora das memórias do projeto (\arq{docs/RESSALVAS.md}, item 11).}}

\section{Os aceleradores}
\label{sec:aceleradores}

A Figura~\ref{fig:fpga-topo} mostra o que há dentro da FPGA: quatro
aceleradores, cada um com as suas memórias e os seus registradores de
controle, todos ligados ao HPS pela interconexão que o Platform Designer gera.

\begin{figure}[htbp]
  \centering
  % Topo da FPGA: cada acelerador com as suas memorias e PIOs.
% Linha por acelerador: [SRAMs de entrada] -> [nucleo] -> [SRAM de saida].
\begin{tikzpicture}[
    ent/.style={mem, minimum width=17mm},
    nuc/.style={caixa, fill=white, minimum width=34mm, minimum height=13mm, font=\scriptsize},
    pio/.style={font=\tiny\ttfamily, text=black!70},
  ]
  \def\xe{0}\def\xn{36mm}\def\xs{72mm}

  % --- convolucao -------------------------------------------------------
  \node[ent, fill=opConv!10] (cx) at (\xe,  3mm) {conv1d\_xn};
  \node[ent, fill=opConv!10] (ch) at (\xe, -3mm) {conv1d\_hn};
  \node[nuc, draw=opConv]    (cc) at (\xn, 0)    {\texttt{conv1d}\\MAC Q15.16\\\(x, h \le 1024 \to y \le 2047\)};
  \node[ent, fill=opConv!10] (cy) at (\xs, 0)    {conv1d\_yn};

  % --- FIR ----------------------------------------------------------------
  \node[ent, fill=opConv!10] (fx) at (\xe, -19mm) {fir\_xn};
  \node[ent, fill=opConv!10] (fh) at (\xe, -25mm) {fir\_hn};
  \node[nuc, draw=opConv]    (fc) at (\xn, -22mm) {\texttt{conv1d} (\texttt{fir\_inst})\\o mesmo RTL da convolução};
  \node[ent, fill=opConv!10] (fy) at (\xs, -22mm) {fir\_yn};

  % --- FFT ----------------------------------------------------------------
  \node[ent, fill=opFft!10] (tx) at (\xe, -41mm) {fft\_xn\_re};
  \node[ent, fill=opFft!10] (ti) at (\xe, -47mm) {fft\_xn\_imag};
  \node[nuc, draw=opFft]    (tc) at (\xn, -44mm) {\texttt{fft\_wrapper} + IP FFT II\\Buffered Burst, N = 1024\\Q15.8, expoente de bloco};
  \node[ent, fill=opFft!10] (ty) at (\xs, -41mm) {fft\_yn\_re};
  \node[ent, fill=opFft!10] (tj) at (\xs, -47mm) {fft\_yn\_imag};

  % --- IIR ----------------------------------------------------------------
  \node[ent, fill=opIir!10] (ix) at (\xe, -63mm) {iir\_xn};
  \node[ent, fill=opIir!10] (ik) at (\xe, -69mm) {iir\_coef};
  \node[nuc, draw=opIir]    (ic) at (\xn, -66mm) {\texttt{iir\_cascade}\\até 16 seções, Q15.16\\1 MAC (15 blocos DSP)};
  \node[ent, fill=opIir!10] (iy) at (\xs, -66mm) {iir\_yn};

  % Setas de dados
  \foreach \a/\b in {cx/cc,ch/cc,fx/fc,fh/fc,tx/tc,ti/tc,ix/ic,ik/ic}
    \draw[seta] (\a.east) -- (\a.east -| \b.west);
  \foreach \a/\b in {cc/cy,fc/fy,ic/iy}
    \draw[seta] (\a.east) -- (\b.west);
  \draw[seta] (tc.east |- ty) -- (ty.west);
  \draw[seta] (tc.east |- tj) -- (tj.west);

  % PIOs de controle (acima de cada nucleo)
  \node[pio, above=0.5mm of cc] {PIO: start, done};
  \node[pio, above=0.5mm of fc] {PIO: start, done, error};
  \node[pio, above=0.5mm of tc] {PIO: start, done, inverse, bfp\_exponent};
  \node[pio, above=0.5mm of ic] {PIO: start, done, error, nsecoes};

  % Interconexao: desce pela esquerda (entradas) e pela direita (saidas)
  \coordinate (bE) at (-12mm, 0);        % barramento da esquerda
  \coordinate (bD) at (\xs+13mm, 0);     % barramento da direita
  \def\ytopo{14mm}
  \draw[thick, cHps] (bE |- 0,-72mm) -- (bE |- 0,\ytopo) -- node[rotulo, above, text=cHps]{interconexão do Platform Designer} (bD |- 0,\ytopo) -- (bD |- 0,-66mm);
  \foreach \y in {3,-3,-19,-25,-41,-47,-63,-69}
    \draw[cHps] (bE |- 0,\y mm) -- (\xe-8.5mm,\y mm);
  \foreach \y in {0,-22,-41,-47,-66}
    \draw[cHps] (bD |- 0,\y mm) -- (\xs+8.5mm,\y mm);

  % HPS, a esquerda, fora da FPGA
  \node[hps, minimum height=40mm, text width=21mm, anchor=east] (hps) at (-24mm,-29mm)
    {HPS\\[2mm]\scriptsize ponte HPS\(\to\)FPGA: memórias\\[2mm] ponte leve: PIOs};
  \draw[<->, very thick, cHps] (hps.east) -- (bE |- hps.east);

  % cantos da FPGA (virgula dentro de fit=... quebraria a lista de opcoes)
  \coordinate (cantoA) at (bE |- 0,\ytopo);
  \coordinate (cantoB) at (bD |- 0,-72mm);
  \begin{scope}[on background layer]
    \node[grupo, draw=cFpga, fill=cFpga!4, fit=(cx)(iy)(ik)(tj)(cantoA)(cantoB),
          inner ysep=5mm,
          label={[font=\scriptsize\bfseries, text=cFpga]above:{FPGA (Cyclone V), relógio de 50 MHz}}] {};
  \end{scope}
\end{tikzpicture}

  \caption{O que há dentro da FPGA: quatro aceleradores, cada um com as suas
    memórias e os seus registradores de controle.}
  \label{fig:fpga-topo}
\end{figure}

\subsection{O protocolo comum}

Todos os aceleradores obedecem ao mesmo combinado com o servidor:

\begin{enumerate}
  \item com \texttt{start} em 0, o servidor escreve as entradas nas memórias;
  \item o servidor leva \texttt{start} a 1; o acelerador detecta a
    \emph{subida} do sinal, zera o \texttt{done} e começa;
  \item o acelerador lê as entradas, calcula, escreve as saídas e leva
    \texttt{done} a 1;
  \item o servidor, que lê \texttt{done} a cada 50\,\textmu s, lê as saídas e devolve
    \texttt{start} a 0.
\end{enumerate}

Como o acelerador reage à subida e não ao nível, um \texttt{start} esquecido em
1 não dispara uma segunda execução. E como o servidor zera todos os
\texttt{start} na primeira requisição depois de a FPGA ser preparada, um valor
que tenha sobrado de antes do bitstream do Morphe também não dispara nada.

Os aceleradores acessam as memórias por pequenos controladores que gastam
alguns ciclos por acesso. É um custo deliberado de simplicidade: numa
convolução, o tempo é dominado por esses acessos, não pela multiplicação. Os
controladores estão em \arq{memory_read_controller.v} e em
\arq{memory_write_controller.v}.

\paragraph{O tempo como sinal de saúde.} Uma convolução de 1024 por 1024
amostras executa cerca de 8,4 milhões de ciclos, uns 168\,ms a 50\,MHz; de ponta
a ponta, com a rede, o cliente mede 214\,ms. Uma resposta muito mais rápida que
isso é sinal de defeito, não de sorte: um bitstream mal compilado já foi visto
levantar o \texttt{done} em 3,3\,ms sem ter calculado nada
(\arq{docs/RESSALVAS.md}, item 13).

\subsection{Convolução e FIR: o \texttt{conv1d}}
\label{sec:acel-conv}

O \arq{conv1d.v} calcula a convolução linear
\[
  y[n] = \sum_{k} x[k]\,h[n-k], \qquad n = 0, 1, \dots, 2046,
\]
com os limites da soma ajustados a cada \(n\) para não ler fora dos vetores. É
uma máquina de estados com um único multiplicador-acumulador: para cada \(n\),
lê \(x[k]\) e \(h[n-k]\), multiplica, acumula e, ao fim da soma, escreve
\(y[n]\).

A aritmética é em Q15.16. O produto de dois números de 32 bits tem 64 bits;
o \texttt{conv1d} guarda os bits 47 a 16, o que realinha o ponto fracionário
\emph{truncando} os 16 bits de baixo. O acumulador tem 32 bits e \emph{não
satura}: quando a soma passa de \(\pm 32768\) ela dá a volta e aparece com o
sinal trocado. Duas consequências práticas:

\begin{itemize}
  \item o erro de truncamento se acumula com o número de termos: numa
    convolução de 1024 por 1024 ele chega a centenas de LSB. É esperado;
  \item com sinais de amplitude \(A\) e \(N\) termos, a saída estoura quando
    \(N A^2 \ge 32768\) --- com \(N = 1024\), a partir de \(A \approx 5{,}7\).
    Uma saída negativa para sinais positivos é esse estouro, não um defeito.
\end{itemize}

Trocar o truncamento por arredondamento reduziria o erro máximo medido de 521
para 25 LSB. A troca foi avaliada e \textbf{deliberadamente não feita}
(14/09/2026): o comportamento numérico do \texttt{conv1d} é o que o TCC
original descreve, e ele fica registrado em \arq{docs/RESSALVAS.md} como
escolha, não como defeito esquecido.

\paragraph{O FIR é o mesmo circuito.} O filtro FIR é uma segunda instância do
\texttt{conv1d} (\texttt{fir\_inst}), com memórias próprias. Filtrar um sinal
por um FIR é convoluí-lo com a resposta ao impulso do filtro, e é exatamente
isso que acontece; a diferença entre as duas operações está no aplicativo, que
oferece para o FIR um projetista de filtros. O projeto ainda contém um IP de FIR
da Intel (\texttt{fir\_ii}) e o seu adaptador, \arq{fir_wrapper.v}, mas a
instância está comentada no \arq{ghrd_top.v}. O registrador
\texttt{fir\_error} sobrou desse adaptador e o \texttt{conv1d} não o aciona.

\subsection{FFT e IFFT: o núcleo da Intel}
\label{sec:acel-fft}

A transformada usa o IP \emph{FFT~II} da Intel, configurado para 1024 pontos
em modo \emph{Buffered Burst}, com entrada de 24 bits e fatores de giro de
18 bits. O \arq{fft_wrapper.v} faz a ponte entre as memórias e o IP: lê as
partes real e imaginária da entrada, entrega-as ao IP amostra por amostra
respeitando o \emph{handshake} do barramento de fluxo, recebe o resultado e o
escreve nas memórias de saída. O registrador \texttt{fft\_inverse} escolhe entre
a transformada direta e a inversa; o resto do caminho é o mesmo.

Duas particularidades desse núcleo aparecem no resultado:

\begin{description}
  \item[Expoente de bloco.] A cada estágio, o IP detecta se os valores
    cresceram e, se preciso, desloca as 1024 amostras \emph{juntas},
    acumulando um expoente único para o bloco (\emph{block floating point}). O
    adaptador captura esse expoente no registrador \texttt{fft\_bfp\_exponent},
    e o servidor o aplica, junto com o fator \(2^{-8}\) do formato Q15.8, antes
    de devolver o resultado em ponto flutuante. É a única conversão numérica
    feita no servidor.
  \item[IFFT sem o fator \(1/N\).] O núcleo entrega a inversa não normalizada,
    \(N\) vezes a IDFT matemática. O cliente aplica o \(1/1024\).
\end{description}

O capítulo~\ref{cap:numerica} explica o que o expoente de bloco e os 24 bits
de entrada significam para a precisão, e por que ela para em torno de 92\,dB.

\subsection{IIR: a cascata de seções de segunda ordem}
\label{sec:acel-iir}

O filtro IIR é a única operação em que a saída depende das saídas anteriores,
e isso muda o que importa no projeto do circuito. Ele é implementado como uma
cascata de até 16 seções de segunda ordem (\emph{biquads}), cada uma em forma
direta~I (Figura~\ref{fig:iir-cascata}):
\[
  y_s[n] = b_0\,x_s[n] + b_1\,x_s[n-1] + b_2\,x_s[n-2]
           - a_1\,y_s[n-1] - a_2\,y_s[n-2],
\]
com a saída de cada seção servindo de entrada para a seguinte. Os coeficientes
vêm da memória \texttt{iir\_coef}, cinco por seção, na ordem
\(b_0, b_1, b_2, a_1, a_2\) (\(a_0 = 1\) fica implícito); o número de seções,
do registrador \texttt{iir\_nsecoes}.

\begin{figure}[htbp]
  \centering
  % Cascata de secoes de 2a ordem e uma secao em forma direta I.
\begin{tikzpicture}[
    sos/.style={caixa, draw=opIir, fill=opIir!10, minimum width=16mm},
    som/.style={circle, draw, inner sep=0pt, minimum size=5mm},
    atr/.style={draw, minimum size=6mm, font=\scriptsize, inner sep=1pt},
    gan/.style={isosceles triangle, draw, inner sep=0.6pt, minimum height=5mm,
                font=\tiny, isosceles triangle apex angle=60},
  ]
  % ---- Cascata -----------------------------------------------------------
  \node (xin) at (0,0) {\(x[n]\)};
  \node[sos, right=8mm of xin] (s1) {seção 1};
  \node[sos, right=8mm of s1]  (s2) {seção 2};
  \node[right=6mm of s2] (dots) {\(\cdots\)};
  \node[sos, right=6mm of dots] (sS) {seção \(S\)};
  \node[right=8mm of sS] (yout) {\(y[n]\)};
  \draw[seta] (xin) -- (s1);
  \draw[seta] (s1) -- (s2);
  \draw[seta] (s2) -- (dots);
  \draw[seta] (dots) -- (sS);
  \draw[seta] (sS) -- (yout);
  \node[rotulo, below=1mm of s2] {cada seção: \(b_0, b_1, b_2, a_1, a_2\) em Q15.16, com \(a_0 = 1\)\\
     até 16 seções, ordem até 32};

  % ---- Uma secao, forma direta I -------------------------------------------
  \begin{scope}[yshift=-26mm, xshift=10mm]
    \node (x) at (0,0) {\(x_s[n]\)};
    \coordinate (xa) at (12mm,0);
    \node[atr] (z1) at (12mm,-11mm) {\(z^{-1}\)};
    \node[atr] (z2) at (12mm,-22mm) {\(z^{-1}\)};
    \node[gan] (b0) at (22mm,0)    {\(b_0\)};
    \node[gan] (b1) at (22mm,-11mm){\(b_1\)};
    \node[gan] (b2) at (22mm,-22mm){\(b_2\)};
    \node[som] (sm) at (38mm,0) {\(+\)};
    \node[caixa, fill=white, font=\scriptsize] (q) at (56mm,0) {arredonda\\satura};
    \coordinate (ya) at (72mm,0);
    \node (y) at (82mm,0) {\(y_s[n]\)};
    \node[atr] (w1) at (72mm,-11mm) {\(z^{-1}\)};
    \node[atr] (w2) at (72mm,-22mm) {\(z^{-1}\)};
    \node[gan, shape border rotate=180] (a1) at (54mm,-11mm) {\(-a_1\)};
    \node[gan, shape border rotate=180] (a2) at (54mm,-22mm) {\(-a_2\)};

    \draw[seta] (x) -- (b0);
    \draw[seta] (xa) -- (z1);  \draw[seta] (z1) -- (z2);
    \draw[seta] (z1) -- (b1);  \draw[seta] (z2) -- (b2);
    \draw[seta] (b0) -- (sm);
    \draw[seta] (b1) -| (sm);
    \draw[seta] (b2) -| ($(sm.south)+(-1.5mm,0)$);
    \draw[seta] (sm) -- (q);
    \draw[seta] (q) -- (y);
    \draw[seta] (ya) -- (w1);  \draw[seta] (w1) -- (w2);
    \draw[seta] (w1) -- (a1);  \draw[seta] (w2) -- (a2);
    \draw[seta] (a1) -| ($(sm.south)+(1.5mm,0)$);
    \draw[seta] (a2) -| ($(sm.south)+(1.5mm,0)$);
    \fill (xa) circle (0.6mm);  \fill (ya) circle (0.6mm);

    \node[rotulo, align=left, anchor=west] at (86mm,-14mm)
      {4 registradores\\de estado por seção;\\as 5 multiplicações\\num só MAC\\compartilhado};
  \end{scope}
\end{tikzpicture}

  \caption{O bloco IIR: cascata de seções de segunda ordem (acima) e a seção em
    forma direta~I (abaixo). Um só multiplicador-acumulador atende todas as
    seções, dois ciclos por seção.}
  \label{fig:iir-cascata}
\end{figure}

Cada escolha desse projeto foi medida antes de ser feita:

\begin{description}
  \item[Seções de segunda ordem, e não um polinômio de ordem \(N\).] A posição
    dos polos de um polinômio de ordem alta é extremamente sensível aos
    coeficientes. Um Butterworth de ordem 22, estável por construção, volta do
    próprio MATLAB com um polo de módulo 1,278 --- instável --- porque achar as
    raízes daquele polinômio já é mal condicionado em ponto flutuante de 64
    bits. A mesma especificação em seções de segunda ordem tem o maior polo em
    0,988.
  \item[Forma direta I.] Tem um acumulador único e largo e arredonda uma vez
    só, na saída da seção. A forma direta II transposta, melhor em ponto
    flutuante, guarda estados já arredondados, e numa realimentação esse erro
    circula.
  \item[Arredondar, e não truncar.] Truncar introduz um viés de meio LSB
    negativo por amostra. Num filtro realimentado esse viés não some: vira um
    deslocamento DC de \(-0{,}5/(1 + a_1 + a_2)\) LSB na saída. No FIR não
    importa, e por isso o \texttt{conv1d} pôde manter o truncamento.
  \item[Q15.16, o mesmo formato do \texttt{conv1d}.] É suficiente: o raio dos
    polos se desloca no máximo cerca de \(10^{-6}\) ao quantizar, mesmo com
    polos em \(|p| = 0{,}9984\).
  \item[Uma amostra atravessa todas as seções antes da próxima.] Assim bastam
    quatro registradores de estado por seção e uma única passagem pela
    memória. O resultado é idêntico, bit a bit, ao de processar seção por
    seção, porque cada seção só vê a saída já arredondada da anterior.
  \item[Um só multiplicador-acumulador, compartilhado.] As cinco multiplicações
    de 32 por 32 bits de uma seção são feitas juntas, e o mesmo circuito
    atende as 16 seções em sequência. O custo em blocos DSP (15, porque cada
    multiplicação de 32 bits ocupa três multiplicadores de 27 bits) não depende
    do número de seções. Cada seção leva dois ciclos: um para selecionar os
    operandos, outro para calcular. Com um ciclo só, o circuito não fechava a
    50\,MHz (49,11\,MHz medidos em 14/09/2026); com dois, fecha a 60,99\,MHz.
\end{description}

Se alguma amostra satura, o registrador \texttt{iir\_error} sobe; o resultado
é devolvido mesmo assim, saturado, e o cliente avisa o aluno. O estado interno é
zerado a cada execução, para que o mesmo sinal dê sempre a mesma saída.

\paragraph{Verificado bit a bit.} O RTL do IIR tem um modelo de referência em
Python que usa inteiros de precisão arbitrária (o acumulador de uma seção
chega a \(2^{65}\), além do que um \texttt{float64} representa exatamente). O
testbench em Icarus Verilog (\arq{Quartus/roda_tb_iir.sh}) e a placa
(\arq{testa_iir_hw.py}) são comparados com esse modelo e precisam coincidir
em todas as amostras: num filtro realimentado, divergir um LSB numa amostra é
divergir para sempre dali em diante, então ``parecido'' não serve de critério.

\section{Periféricos ainda não usados: ADC e codec}
\label{sec:adc}

A placa tem dois conversores que o Morphe ainda não usa. Esta seção descreve o
estado de hoje e o que falta; a Figura~\ref{fig:adc} resume as duas coisas.

\begin{figure}[htbp]
  \centering
  % ADC: o que existe hoje (linha cheia) e o que falta (tracejado).
\begin{tikzpicture}[node distance=6mm]
  % ---- Hoje ---------------------------------------------------------------
  \node[caixa, fill=white, text width=16mm] (con) {conector\\analógico\\\scriptsize CH0--CH7};
  \node[caixa, fill=black!6, right=of con, text width=23mm] (ltc) {LTC2308\\\scriptsize 12 bits, 8 canais\\0--4,095\,V};
  \node[fpga, right=18mm of ltc, text width=31mm] (up) {controlador do\\University Program\\\scriptsize\itshape instância comentada};

  \draw[seta] (con) -- (ltc);
  \draw[<->, thick] (ltc) -- node[rotulo, above]{SPI, 4 fios}
                              node[rotulo, below]{\texttt{CONVST}\\\texttt{SCLK DIN}\\\texttt{DOUT}} (up);

  % ---- O que falta ----------------------------------------------------------
  \node[futuro, right=8mm of up, text width=18mm] (pio) {PIO\\\scriptsize leitura avulsa};
  \node[futuro, below=14mm of ltc, text width=24mm] (ctl) {controlador próprio\\\scriptsize CONVST a \(1/f_s\)};
  \node[futuro, text width=18mm] (ram) at (ctl -| up) {RAM de\\captura};
  \node[futuro, text width=18mm] (op) at (ctl -| pio) {opcodes;\\janela\\``Aquisição''};

  \draw[seta, dashed, cFuturo] (up) -- (pio);
  \draw[seta, dashed, cFuturo] (pio) -- (op);
  \draw[seta, dashed, cFuturo] (ltc) -- (ctl);
  \draw[seta, dashed, cFuturo] (ctl) -- (ram);
  \draw[seta, dashed, cFuturo] (ram) -- (op);

  \node[rotulo, anchor=north, text=cFuturo!70!black, align=left] at ($(ctl.south)!0.5!(op.south)+(0,-2mm)$)
    {etapa 1: leitura avulsa, o \emph{voltímetro} (linha de cima)\\
     etapa 2: captura a \(f_s\) fixa em RAM (linha de baixo)};
\end{tikzpicture}

  \caption{O ADC da placa: o que existe (linha cheia) e o que falta para
    capturar sinais a uma taxa fixa (tracejado).}
  \label{fig:adc}
\end{figure}

\paragraph{O ADC.} O conversor analógico-digital da DE1-SoC é o LTC2308: 12
bits, 8 canais, por aproximações sucessivas, ligado à FPGA por uma interface
serial de quatro fios (\texttt{ADC\_CONVST}, \texttt{ADC\_SCLK},
\texttt{ADC\_DIN}, \texttt{ADC\_DOUT}). As entradas analógicas chegam por um
conector próprio da placa. O projeto contém o controlador do \emph{University
Program} da Intel para esse conversor, configurado para entrada simples e
unipolar: com a referência interna, a faixa é de 0 a 4,095\,V, 1\,mV por passo.
\textbf{Tensão negativa não é lida e pode danificar a entrada}; um sinal
alternado precisa de um deslocamento DC antes de chegar ao conector.

Hoje a instância desse controlador está comentada no \arq{ghrd_top.v}: o ADC
não faz parte do bitstream. E mesmo ativo, ele não serviria para amostrar um
sinal: converte sem parar, sozinho, e entrega apenas o último valor lido de
cada canal, sem controle do instante da conversão. Uma incerteza de alguns
microssegundos no instante da amostra limita a medida de um tom de 5\,kHz a
cerca de 25\,dB. Capturar a uma taxa de amostragem fixa exige um controlador
próprio, que dispare a conversão a cada \(1/f_s\) e guarde as amostras numa
memória da FPGA. O plano, em duas etapas --- uma leitura avulsa, como um
voltímetro, e depois a captura temporizada --- está em \arq{docs/ADC.md}.

\paragraph{O codec de áudio.} A placa tem também um codec de áudio (o WM8731,
segundo a documentação da DE1-SoC), com entrada e saída de linha, ligado à FPGA
por uma interface de áudio serial e configurado por I\textsuperscript{2}C. Os
pinos estão declarados no \arq{ghrd_top.v} e mais nada: nenhum sinal do codec é
usado pelo projeto. Seria o caminho natural para uma saída de áudio --- ouvir o
sinal filtrado --- e fica como trabalho futuro.


% =============================================================================
\chapter{O servidor na placa}
\label{cap:servidor}
% =============================================================================

\begin{previsto}
O \arq{morphe_server.c}: mapeamento de \arq{/dev/mem}, laço
\emph{accept}--atende--fecha, um cliente por vez, fila do sistema operacional
(\arq{listen(4)}). Validação do cabeçalho, tempo-limite de 5\,s por operação.
A marca \arq{/var/run/morphe-fpga-preparada} e por que o servidor recusa
operar sem ela (o travamento do barramento com o bitstream de fábrica,
21/09). Compilação na própria placa e \emph{autostart}. Tabela dos códigos de
status.
\end{previsto}

% =============================================================================
\chapter{O protocolo Morphe-TCP}
\label{cap:protocolo}
% =============================================================================

\begin{previsto}
Uma conexão por operação. Cabeçalhos de 20 bytes em \emph{big-endian}
(Figura~\ref{fig:protocolo}), tabela de opcodes (CONV, FFT, PING, FIR, IFFT,
IIR) com o formato do \emph{payload} de cada um, tabela de status. Como
acrescentar um opcode: os quatro lugares que mudam juntos
(\arq{morphe_protocol.h}, \arq{morphe_server.c}, \arq{morphe_protocol.py},
\arq{morphe_config.*}).
\end{previsto}

\begin{figure}[htbp]
  \centering
  % Cabecalhos do protocolo Morphe-TCP (C/morphe_protocol.h). 1 byte = 6 mm.
\begin{tikzpicture}[
    campo/.style={draw, minimum height=9mm, inner sep=0pt, align=center,
                  font=\scriptsize},
  ]
  \def\u{6.2mm}
  % campo(inicio, bytes, texto, cor)
  \newcommand{\campo}[5]{%
    \node[campo, fill=#4, minimum width=#2*\u, anchor=west] at (#1*\u, #5) {#3};}

  % ---- Requisicao ------------------------------------------------------
  \node[anchor=east, font=\small\bfseries] at (-2mm,0) {requisição};
  \campo{0}{4}{\texttt{magic}\\\texttt{"MRPM"}}{cCliente!12}{0}
  \campo{4}{2}{\texttt{version}\\1}{white}{0}
  \campo{6}{2}{\texttt{opcode}}{opIir!15}{0}
  \campo{8}{2}{\texttt{dtype}}{white}{0}
  \campo{10}{2}{\texttt{flags}}{white}{0}
  \campo{12}{4}{\texttt{n\_x}\\amostras de \(x\)}{white}{0}
  \campo{16}{4}{\texttt{n\_h}\\\(h\) ou seções}{white}{0}

  % ---- Resposta --------------------------------------------------------
  \node[anchor=east, font=\small\bfseries] at (-2mm,-26mm) {resposta};
  \campo{0}{4}{\texttt{magic}\\\texttt{"MRPN"}}{cHps!12}{-26mm}
  \campo{4}{2}{\texttt{version}}{white}{-26mm}
  \campo{6}{2}{\texttt{opcode}}{opIir!15}{-26mm}
  \campo{8}{2}{\texttt{dtype}}{white}{-26mm}
  \campo{10}{2}{\texttt{status}}{opConv!15}{-26mm}
  \campo{12}{4}{\texttt{n\_out}}{white}{-26mm}
  \campo{16}{4}{\texttt{extra}\\saturou (IIR)}{white}{-26mm}

  % Regua de bytes
  \foreach \b in {0,4,8,12,16,20}
    \node[font=\tiny, text=black!60] at (\b*\u, 7mm) {\b};
  \draw[black!40] (0,5.5mm) -- (20*\u,5.5mm);

  % Payload, logo depois de cada cabecalho
  \node[campo, fill=black!5, dashed, minimum width=20*\u, anchor=west] at (0,-9mm)
    {\emph{payload}: \(x\) e \(h\), espectro complexo ou coeficientes --- int32, 4 bytes por valor};
  \node[campo, fill=black!5, dashed, minimum width=20*\u, anchor=west] at (0,-35mm)
    {\emph{payload}: \(y\) (int32 ou float32), ou a mensagem de erro se \texttt{status} \(\ne 0\)};
\end{tikzpicture}

  \caption{Os cabeçalhos de requisição e de resposta, em bytes.}
  \label{fig:protocolo}
\end{figure}

% =============================================================================
\chapter{Representação numérica}
\label{cap:numerica}
% =============================================================================

\begin{previsto}
Os dois formatos de ponto fixo (Figura~\ref{fig:ponto-fixo}): Q15.8 na FFT e
na IFFT, Q15.16 na convolução, no FIR e no IIR. A regra dos 6,02\,dB por bit,
a normalização da entrada e o ponto fixo em bloco do IP da FFT; por que o
teto é \(\sim\)92\,dB. A IFFT sem o fator \(1/N\). Tabela dos erros medidos e
regras práticas. Resume \arq{docs/PRECISAO-NUMERICA.md}.
\end{previsto}

\begin{figure}[htbp]
  \centering
  % Q15.8 dentro de um int32 (FFT/IFFT) e Q15.16 (conv, FIR, IIR). 1 bit = 3,6 mm.
\begin{tikzpicture}
  \def\u{3.4mm}
  % faixa(bit_inicial_da_esquerda, n_bits, texto, cor, y)
  \newcommand{\faixa}[5]{%
    \draw[fill=#4] (#1*\u,#5) rectangle ++(#2*\u,6mm);
    \foreach \k in {1,...,#2} \draw[black!25] ({(#1+\k)*\u},#5) -- ++(0,6mm);
    \node[font=\scriptsize] at ({(#1+#2/2)*\u},#5+3mm) {#3};}

  % ---- Q15.8 --------------------------------------------------------------
  \node[anchor=east, font=\small, align=right] at (-2mm,3mm) {\textbf{Q15.8}\\\scriptsize FFT, IFFT};
  \faixa{0}{8}{extensão de sinal}{black!8}{0}
  \faixa{8}{1}{s}{opConv!25}{0}
  \faixa{9}{15}{parte inteira (15)}{cCliente!15}{0}
  \faixa{24}{8}{fração (8)}{cHps!18}{0}
  \draw[decorate, decoration={brace, mirror, amplitude=4pt}]
    (8*\u,-1mm) -- node[below=4pt, font=\scriptsize]{24 bits que o IP da FFT lê (\texttt{sink\_real[23:0]})} (32*\u,-1mm);
  \node[anchor=west, font=\scriptsize, align=left] at (33*\u,3mm) {passo \(2^{-8}\)\\faixa \(\pm 32768\)};

  % ---- Q15.16 -------------------------------------------------------------
  \node[anchor=east, font=\small, align=right] at (-2mm,-17mm) {\textbf{Q15.16}\\\scriptsize conv, FIR, IIR};
  \faixa{0}{1}{s}{opConv!25}{-20mm}
  \faixa{1}{15}{parte inteira (15)}{cCliente!15}{-20mm}
  \faixa{16}{16}{fração (16)}{cHps!18}{-20mm}
  \node[anchor=west, font=\scriptsize, align=left] at (33*\u,-17mm) {passo \(2^{-16}\)\\faixa \(\pm 32768\)};

  % regua
  \foreach \b/\p in {31/0,24/7,16/15,8/23,0/31}
    \node[font=\tiny, text=black!60] at ({(\p+0.5)*\u},8mm) {\b};
\end{tikzpicture}

  \caption{Os dois formatos de ponto fixo, bit a bit.}
  \label{fig:ponto-fixo}
\end{figure}

% =============================================================================
\chapter{O aplicativo cliente}
\label{cap:cliente}
% =============================================================================

\section{Janelas e módulos}
\begin{previsto}
O \arq{morphe_app.py} e as janelas (gerador de sinais, convolução, FFT, IFFT,
FIR com projetista, IIR com projetista, comparador), com uma captura de tela
de cada --- as únicas figuras que não são desenhadas em \LaTeX. De onde vem o
sinal: gerador, digitado ou arquivo (\arq{.csv}, \arq{.txt}, \arq{.npy},
\arq{.wav}). Os módulos por trás das janelas (\arq{dsp_core}, \arq{blocos},
\arq{fir_design}, \arq{iir_design}, \arq{morphe_export}). O pacote
\arq{.mrph} e o comparador.
\end{previsto}

\section{Sinais maiores que o hardware: processamento em blocos}
\begin{previsto}
O hardware processa 1024 amostras por vez; o cliente divide e recompõe. Três
técnicas, uma por tipo de operação, cada uma com a sua figura e o seu
argumento de exatidão: \emph{overlap-add} para convolução e FIR; algoritmo de
quatro passos para FFT e IFFT (a mesma DFT, não uma aproximação); blocos com
aquecimento para o IIR (bit a bit igual ao modelo nos mesmos blocos).
\end{previsto}

\begin{figure}[htbp]
  \centering
  % Overlap-add: x em blocos de L, cada um convoluido na placa com h (M amostras).
\begin{tikzpicture}[
    blk/.style={draw, minimum height=6mm, inner sep=0pt, font=\scriptsize},
  ]
  \def\L{30mm}\def\T{10mm} % comprimento do bloco e da cauda (M-1) no desenho

  % Sinal de entrada
  \node[anchor=east, font=\small] at (-2mm,0) {\(x[n]\)};
  \foreach \i/\c in {0/opConv, 1/opFft, 2/opIir} {
    \node[blk, fill=\c!15, minimum width=\L, anchor=west] at (\i*\L,0) {bloco \the\numexpr\i+1\relax: \(L\) amostras};
  }
  \node[font=\small] at (3*\L+6mm,0) {\(\cdots\)};

  % Convolucoes na placa (deslocadas)
  \foreach \i/\c in {0/opConv, 1/opFft, 2/opIir} {
    \pgfmathsetmacro{\y}{-12-9*\i}
    \node[blk, fill=\c!15, minimum width=\L, anchor=west] (c\i) at (\i*\L,\y mm) {\(x_{\the\numexpr\i+1\relax} * h\)};
    \node[blk, fill=\c!35, minimum width=\T, anchor=west] at (c\i.east) {cauda};
    \node[font=\scriptsize, anchor=east, text=black!60] at (-2mm,\y mm) {na placa};
  }

  % Soma
  \draw[thick] (0,-42mm) -- (3*\L+\T,-42mm);
  \node[anchor=east, font=\small] at (-2mm,-48mm) {\(y[n]\)};
  \node[blk, fill=black!5, minimum width=3*\L+\T, anchor=west] at (0,-48mm)
    {soma: onde duas caudas se sobrepõem, as amostras se somam};

  % Marcas de sobreposicao
  \foreach \i in {1,2}
    \draw[densely dashed, black!50] (\i*\L,-8mm) -- (\i*\L,-45mm);
  \draw[decorate, decoration={brace, amplitude=3pt}]
    (\L,-37mm) -- node[below=3pt, font=\scriptsize]{\(M-1\)} (\L+\T,-37mm);
\end{tikzpicture}

  \caption{\emph{Overlap-add}: cada bloco de \(L\) amostras é convoluído na
    placa; as caudas de \(M-1\) amostras se somam.}
  \label{fig:overlap-add}
\end{figure}

\begin{figure}[htbp]
  \centering
  % FFT longa pelo algoritmo de quatro passos (Python/blocos.py, fft_longa).
% N = M * 1024; n = n1 + M*n2; k = 1024*k1 + k2.
\begin{tikzpicture}[
    passo/.style={font=\scriptsize\bfseries, align=center},
    lin/.style={draw, minimum width=34mm, minimum height=5mm, inner sep=0pt, font=\tiny},
  ]
  % ---- 1: subsequencias -> M FFTs de 1024 na placa --------------------------
  \node[passo] at (17mm,8mm) {1. \(M\) subsequências\\\(x[n_1 + M n_2]\)};
  \foreach \r in {0,1,2} {
    \node[lin, fill=opFft!12] (a\r) at (17mm,-\r*6mm) {\(n_1 = \r\): FFT na placa};
  }
  \node at (17mm,-18mm) {\(\vdots\)};
  \node[lin, fill=opFft!12] (a3) at (17mm,-24mm) {\(n_1 = M-1\)};

  % ---- 2: giro ---------------------------------------------------------------
  \node[passo] at (58mm,8mm) {2. multiplica por\\\(W_N^{\,n_1 k_2}\)};
  \node[caixa, fill=white, minimum width=20mm, minimum height=30mm] (g) at (58mm,-12mm)
    {fatores\\de giro\\\scriptsize (no PC)};

  % ---- 3: DFTs de M pontos por coluna ------------------------------------------
  \node[passo] at (100mm,8mm) {3. DFT de \(M\) pontos\\em cada coluna \(k_2\)};
  \draw[fill=cCliente!8] (83mm,-26.5mm) rectangle (117mm,2.5mm);
  \foreach \c in {0,...,8}
    \draw[cCliente, thick, ->] ({85mm+\c*3.8mm},1.5mm) -- ({85mm+\c*3.8mm},-25.5mm);
  \node[font=\tiny, fill=cCliente!8] at (100mm,-12mm) {1024 colunas, no PC};

  % ---- 4: leitura por linhas -----------------------------------------------------
  \node[passo] at (140mm,8mm) {4. lê por linhas:\\\(X[1024 k_1 + k_2]\)};
  \draw[fill=black!4] (125mm,-26.5mm) rectangle (155mm,2.5mm);
  \foreach \r in {0,...,4}
    \draw[thick, ->] (126.5mm,{0.5mm-\r*6mm}) -- (153.5mm,{0.5mm-\r*6mm});

  \draw[seta] (34.5mm,-12mm) -- (g.west);
  \draw[seta] (g.east) -- (83mm,-12mm);
  \draw[seta] (117mm,-12mm) -- (125mm,-12mm);

  \node[rotulo, text width=150mm, align=left, anchor=north west] at (0,-31mm)
    {Não é aproximação: é a DFT de \(N\) pontos com a conta dividida
     (Cooley--Tukey com \(N_1 = M\), \(N_2 = 1024\)). A IFFT longa usa a mesma
     decomposição, com os fatores conjugados e \(1/N = 1/1024 \cdot 1/M\).};
\end{tikzpicture}

  \caption{FFT de \(N = M\cdot 1024\) pontos pelo algoritmo de quatro passos:
    \(M\) FFTs de 1024 na placa, fatores de giro e DFTs de \(M\) pontos no PC.}
  \label{fig:quatro-passos}
\end{figure}

\begin{figure}[htbp]
  \centering
  % IIR por blocos com aquecimento (Python/blocos.py, iir_por_blocos).
\begin{tikzpicture}[
    blk/.style={draw, minimum height=6mm, inner sep=0pt, font=\scriptsize, anchor=west},
  ]
  \def\B{36mm}\def\A{10mm} % bloco da placa (1024) e aquecimento no desenho

  \node[anchor=east, font=\small] at (-2mm,0) {\(x[n]\)};
  \draw[fill=black!4] (0,-3mm) rectangle (3*\B-2*\A+8mm,3mm);
  \node[font=\scriptsize] at (1.5*\B,0) {sinal de arquivo, qualquer tamanho};

  % bloco 1: sem aquecimento
  \node[blk, fill=opIir!25, minimum width=\B] at (0,-10mm) {bloco 1: 1024 guardadas};
  % bloco 2
  \node[blk, fill=black!10, minimum width=\A] at (\B-\A,-17mm) {};
  \node[blk, fill=opIir!25, minimum width=\B-\A] at (\B,-17mm) {bloco 2: \(1024 - A\)};
  % bloco 3
  \node[blk, fill=black!10, minimum width=\A] at (2*\B-2*\A,-24mm) {};
  \node[blk, fill=opIir!25, minimum width=\B-\A] at (2*\B-\A,-24mm) {bloco 3};

  \draw[decorate, decoration={brace, mirror, amplitude=3pt}]
    (\B-\A,-20.5mm) -- node[below=3pt, font=\scriptsize, align=center]{aquecimento \(A\):\\descartado} (\B,-20.5mm);

  \node[rotulo, anchor=west, align=left, text width=48mm] at (3*\B-2*\A+12mm,-17mm)
    {\(A\) sai do polo de maior módulo: o transitório de estado zero cai abaixo de
     meio LSB do Q15.16 antes da primeira amostra guardada.};
\end{tikzpicture}

  \caption{IIR por blocos: cada bloco começa algumas amostras antes
    (\emph{aquecimento}) e essa parte da saída é descartada.}
  \label{fig:iir-blocos}
\end{figure}

% =============================================================================
\chapter{Várias placas, vários alunos}
\label{cap:placas}
% =============================================================================

\begin{previsto}
Descoberta (lista em \arq{.morphe-estado/placas}, varredura como último
recurso), escolha pela sonda mais rápida
(Figura~\ref{fig:escolha-placa}), a fila do servidor e o que ela faz com a
latência. Tabela das medições de 22/09 (6 clientes, 90/90 respostas
conferidas). O que não existe: alocação de placas, concorrência no servidor,
descoberta sem varredura. Resume \arq{docs/GERENCIAMENTO-PLACAS.md}.
\end{previsto}

\begin{figure}[htbp]
  \centering
  % Escolha da placa pela sonda mais rapida (Python/tcp_panel.py).
\begin{tikzpicture}[
    ev/.style={font=\scriptsize, align=center},
  ]
  % eixo de tempo
  \draw[->, black!60] (0,-26mm) -- (125mm,-26mm) node[below left, font=\scriptsize]{tempo};
  \node[cliente, text width=18mm] (cl) at (-14mm,-9mm) {cliente\\\scriptsize \texttt{OP\_PING}\\em paralelo};

  % Placa 1: ocupada com uma convolucao
  \node[anchor=east, font=\scriptsize\bfseries] at (-2mm,0) {};
  \node[hps, anchor=west, minimum width=10mm] at (0,0) {placa 1};
  \draw[fill=opConv!20] (14mm,-3mm) rectangle (100mm,3mm);
  \node[ev] at (57mm,0) {ocupada: convolução de outro aluno (\(\approx\)214\,ms)};
  \draw[fill=cHps!30] (100mm,-3mm) rectangle (106mm,3mm);
  \node[ev, anchor=west] at (107mm,0) {PING};

  % Placa 2: livre
  \node[hps, anchor=west, minimum width=10mm] at (0,-16mm) {placa 2};
  \draw[fill=cHps!30] (16mm,-19mm) rectangle (22mm,-13mm);
  \node[ev, anchor=west] at (23mm,-16mm) {PING responde em milissegundos \(\Rightarrow\) \textbf{escolhida}};

  \draw[seta, cRede] (cl.east) -- (0,0);
  \draw[seta, cRede] (cl.east) -- (0,-16mm);
\end{tikzpicture}

  \caption{Escolha da placa: o cliente sonda todas em paralelo e fica com a
    que responder primeiro. Uma placa ocupada só responde ao fim da operação
    em curso.}
  \label{fig:escolha-placa}
\end{figure}

% =============================================================================
\chapter{Preparação e operação do laboratório}
\label{cap:laboratorio}
% =============================================================================

\begin{previsto}
O \arq{morphe-up.sh} passo a passo (Figura~\ref{fig:morphe-up}) e as suas
opções (\opt{board}, \opt{skip-fpga}, \opt{down},
\opt{setup-ssh}); o \emph{tether} da licença, um por cabo e por conta. A
instalação de turma em \arq{/opt/morphe} (conta do coordenador e conta dos
alunos). O \arq{instala-quartus.sh} e o Quartus~20.1 em todos os PCs.
Incorporar uma placa nova (\arq{docs/NOVA-PLACA.md}). Tabela sintoma
\(\to\) causa \(\to\) o que fazer.
\end{previsto}

\begin{figure}[htbp]
  \centering
  % Os passos do morphe-up.sh, na ordem em que o script os imprime.
\begin{tikzpicture}[
    p/.style={caixa, fill=white, text width=54mm, font=\scriptsize},
    onde/.style={font=\tiny\bfseries, anchor=east},
    node distance=4mm,
  ]
  \node[p, draw=cCliente] (p0) {confere a conta (recusa \texttt{root}) e acha o Quartus 20.1};
  \node[p, draw=cFpga, below=of p0] (p1) {programa a FPGA pelo JTAG\\(\texttt{quartus\_pgm} + \emph{tether} do cabo)};
  \node[p, draw=cRede, below=of p1] (p2) {confere o caminho até a placa (rede, SSH)};
  \node[p, draw=cHps, below=of p2] (p3) {envia e compila o servidor na placa\\(\texttt{make -B})};
  \node[p, draw=cHps, below=of p3] (p4) {(re)inicia o servidor e grava a marca\\\texttt{/var/run/morphe-fpga-preparada}};
  \node[p, draw=cCliente, below=of p4] (p5) {verifica a plataforma (\texttt{morphe\_ping})\\e registra a placa em \texttt{.morphe-estado/placas}};

  \foreach \a/\b in {p0/p1,p1/p2,p2/p3,p3/p4,p4/p5} \draw[seta] (\a) -- (\b);

  \node[onde, text=cCliente] at ($(p0.west)+(-2mm,0)$) {estação};
  \node[onde, text=cFpga]    at ($(p1.west)+(-2mm,0)$) {JTAG};
  \node[onde, text=cRede]    at ($(p2.west)+(-2mm,0)$) {rede};
  \node[onde, text=cHps]     at ($(p3.west)+(-2mm,0)$) {placa (SSH)};
  \node[onde, text=cHps]     at ($(p4.west)+(-2mm,0)$) {placa (SSH)};
  \node[onde, text=cCliente] at ($(p5.west)+(-2mm,0)$) {estação};

  \node[rotulo, anchor=west, text width=48mm, align=left] at ($(p1.east)+(6mm,0)$)
    {O \emph{tether} precisa ficar vivo: sem ele o bitstream de avaliação para
     depois de uma hora --- a FFT e, com ela, todas as operações.};
  \node[rotulo, anchor=west, text width=48mm, align=left] at ($(p4.east)+(6mm,0)$)
    {Sem a marca o servidor recusa toda operação (status 8): a placa recém-ligada
     ainda está com o bitstream de fábrica.};
\end{tikzpicture}

  \caption{O que o \arq{morphe-up.sh} faz, em ordem.}
  \label{fig:morphe-up}
\end{figure}

% =============================================================================
\chapter{Desenvolvimento e verificação}
\label{cap:desenvolvimento}
% =============================================================================

\begin{previsto}
Para quem vai estender o sistema. Recompilar o bitstream (Platform Designer,
Quartus, relatório de tempo; resume \arq{docs/COMPILAR-IIR.md}). Testbenches
em Icarus Verilog com vetores gerados em Python. Os testes contra a placa
(\arq{testa_iir_hw}, \arq{testa_blocos}, \arq{testa_concorrencia},
\arq{testa_ifft_roundtrip}, \arq{testa_bit_inverse}). A
\arq{PROVENIENCIA.sha256}: como provar que o que está na placa saiu do
repositório. Fluxo de trabalho com o git (push do Windows, \emph{pull} na
estação).
\end{previsto}

% =============================================================================
\chapter{Limitações e trabalho futuro}
% =============================================================================

\begin{previsto}
O que o sistema não faz hoje, cada item com o motivo: alocação de placas;
aquisição pelo ADC (as etapas planejadas); saída de áudio pelo codec; licença
de avaliação da FFT (o \emph{tether}); tamanho fixo de 1024 no hardware.
\end{previsto}

% =============================================================================
\appendix
\chapter{Estrutura do repositório}
\label{ap:repositorio}
% =============================================================================

\begin{previsto}
A árvore abaixo, com uma linha por arquivo importante. A regra de que os
\arq{iir_*.v} não saem de \arq{Quartus/}.
\end{previsto}

% Arvore do repositorio (estado de de9bc41). Pacote dirtree.
\renewcommand*\DTstylecomment{\small\rmfamily\itshape\color{black!65}}
\renewcommand*\DTstyle{\small\ttfamily}
\dirtree{%
.1 morphe/.
.2 README.md\DTcomment{ponto de entrada}.
.2 morphe-up.sh\DTcomment{prepara a placa: FPGA, servidor, marca}.
.2 instala-quartus.sh\DTcomment{Quartus 20.1 para todas as contas de um PC}.
.2 deploy.sh, gera\_proveniencia.sh, PROVENIENCIA.sha256.
.2 Python/\DTcomment{o aplicativo}.
.3 morphe\_app.py\DTcomment{janela principal}.
.3 *\_window.py\DTcomment{uma janela por operação}.
.3 morphe\_protocol.py\DTcomment{cliente TCP}.
.3 dsp\_core.py, blocos.py\DTcomment{ponto fixo, referência, sinais longos}.
.3 fir\_design.py, iir\_design.py\DTcomment{projetistas de filtros}.
.3 ferramentas/\DTcomment{testes contra a placa, figuras, vetores}.
.2 C/\DTcomment{o servidor da placa}.
.3 morphe\_server.c, morphe\_protocol.h, morphe\_config.h.
.3 hps\_0.h\DTcomment{gerado do Platform Designer}.
.3 autostart/.
.2 Quartus/\DTcomment{o hardware}.
.3 ghrd\_top.v\DTcomment{topo: instancia os aceleradores}.
.3 conv1d.v, fft\_wrapper.v, iir\_*.v\DTcomment{os aceleradores}.
.3 soc\_system.qsys, soc\_system.qsf.
.3 tb\_iir\_*.v, roda\_tb\_iir.sh\DTcomment{testbenches}.
.3 output\_files/\DTcomment{o bitstream}.
.2 exemplos/\DTcomment{pacotes .mrph e sinais de exemplo}.
.2 docs/\DTcomment{documentação técnica e diário}.
.2 legado/\DTcomment{do projeto original, fora do produto}.
}


\chapter{Referência rápida}
\begin{previsto}
Os comandos do dia a dia numa página: preparar a placa, abrir o aplicativo,
conferir um clone, rodar os testes, desligar a plataforma.
\end{previsto}

\end{document}
